%==============
\begin{center}
{\LARGE \bf Computational Physics
}
\end{center}
%==============
\bigskip

{\bf Organizer}

\begin{itemize}

\item
V.G. Ganzha (ganzha@hrz.uni-kassel.de)

\item
Claudio Maccone (cmaccone@to.alespazio.it)

\item
Rolf Mertig (rolfm@nikhefh.nikhef.nl) 

\item
E.V. Vorozhtsov (adm@itam.nsk.su)

\end{itemize}

\medskip

\bigskip

{\bf Description}
\medskip

Computer algebra techniques/softwares have been used in theoretical/computational physics (such
as general relativity, quantum physics, solid state physics, etc) since the late sixties. Nowadays, due
to dramatic improvements in hardware, algorithms and software systems, more and more
non-trivial results, e.g. in perturbation theory and stability theory in PDE, are achieved by using
computer algebraic methods/programs. While for physicists the results and the interpretation of
their research is naturally most important, it is realized at a larger scale now that computer algebra
systems and special-purpose programs written in the corresponding mathematical high-level
languages can be extremely useful (and fun) when doing calculationally intensive investigations.
This session has as its goal to bring together researchers who make serious use of computer algebra
theory/systems in theorical or computational physics. 

\bigskip
{\bf Talks}
\medskip

Date: July 18th (Thursday) 


\begin{description}
\item[08:30-08:55] \ \\
  Claudio Maccone\\
  {\bf Computer Algebra and General Relativity for Future Spaceflight Investigatio.} 
\item[08:55-09:20] \ \\
  Rolf Mertig
\\
  {\bf FeynCalc 3.0 - A Mathematica Package for Feynman Diagram Calculations in High
    Energy Physics.} 
\item[09:20-09:45] \ \\
  Michael Trott
\\
  {\bf High order WKB approximation and singular perturbation theory.}
  
\item[09:45-10:20] \ \\
  \ \\
  {\bf Coffee Break}
\item[10:20-10:45] \ \\
  Victor G. Ganzha, K. Schaub, Evgenii V. Vorozhtsov
\\
  {\bf The Use of Mathematica for Computation of Stability Regions with Guaranteed
    Accuracy.} 
\item[10:45-11:10 ] \ \\
  Victor G. Ganzha, Evgenii V. Vorozhtsov \\
  {\bf Curvilinear Grid Topology Effect on the Stability of Difference Schemes.} 
\item[11:10-11:35] \ \\
  Hans-Joachim Bungartz\\
  {\bf Rapid Prototyping for the Construction of Higher Order Finite Element Methods on
    Sparse Grids.} 
\item[11:35-12:00] \ \\
  Jean-Antoine D\'esid\'eri, Margarita Spiridonova\\
  {\bf Symbolic Computation of Linear and Nonlinear Modified (Partial Differential)
    Equations.} 
\end{description}

